\documentclass[11pt]{article}
\usepackage[utf8]{inputenc}
\usepackage[T1]{fontenc}
\usepackage{amsmath,amssymb}
\usepackage{graphicx}
\usepackage{booktabs}
\usepackage{hyperref}
\usepackage[margin=1in]{geometry}
\usepackage{natbib}
\usepackage{authblk}

\title{Disseminating Cells in Human Tumours Acquire an EMT Stem Cell State That Predicts Metastasis}

\author[1]{Author A}
\author[1]{Author B}
\author[1,2]{Author C}
\affil[1]{Department of Oncology, Medical University of Graz, Graz, Austria}
\affil[2]{Institute of Pathology, Medical University of Graz, Graz, Austria}

\date{}

\begin{document}
\maketitle

\begin{abstract}
Epithelial--mesenchymal transition (EMT) and cancer stem cell (CSC) properties are thought to drive metastasis, but direct observation of disseminating EMT CSCs in human tumour specimens has been lacking. Here, we used multi-marker immunofluorescence staining for EpCAM, CD24, and Vimentin to characterize over 12,000 imaging fields from 84 human oral cancer specimens. We identified triple-positive EpCAM$^+$/CD24$^+$/Vimentin$^+$ cells at the invasive front that exhibit a hybrid epithelial--mesenchymal phenotype consistent with an EMT stem cell state. These disseminating cells were significantly enriched in metastatic compared to non-metastatic specimens. An artificial neural network trained on disseminating cell counts and marker co-expression features predicted metastatic outcome with 87--89\% cross-validated accuracy. Our findings provide the first direct observation of EMT CSCs actively disseminating from human tumour tissue and demonstrate that their abundance has clinical predictive value for metastasis.
\end{abstract}

\section{Introduction}

Metastasis is the primary cause of cancer-related mortality, yet the cellular mechanisms driving tumour cell dissemination from primary sites remain incompletely understood in human cancers \citep{lambert2017emerging}. The epithelial--mesenchymal transition (EMT) has emerged as a central mechanism by which epithelial cancer cells acquire the motility and invasive properties necessary for dissemination \citep{nieto2016emt}. Concurrently, the cancer stem cell (CSC) model posits that a subpopulation of tumour cells with self-renewal capacity drives tumour initiation, therapy resistance, and metastasis \citep{batlle2017cancer}.

A growing body of evidence suggests that EMT and CSC properties are intertwined: cells undergoing EMT acquire stem cell-like features, and CSCs frequently exhibit mesenchymal characteristics \citep{mani2008mesenchyme}. However, the relationship between EMT, stemness, and metastatic dissemination has been established primarily in experimental models. Whether EMT CSCs can be directly identified at the invasive front of human tumours, where dissemination occurs, and whether their presence predicts clinical metastatic outcome, has remained an open question.

In oral squamous cell carcinoma (OSCC), regional lymph node metastasis is a major prognostic factor, yet reliable biomarkers for predicting metastasis at the time of primary tumour diagnosis are lacking \citep{leemans2011biology}. Identifying the cells that are actively disseminating from the tumour could provide both biological insight into the metastatic process and clinically useful biomarkers.

Here, we develop a multi-marker immunofluorescence approach to identify and quantify disseminating cells at the invasive front of human oral cancers. Using the combination of EpCAM (epithelial), CD24 (stem cell), and Vimentin (mesenchymal), we identify triple-positive cells with a hybrid EMT CSC phenotype. We demonstrate that these cells are enriched in metastatic specimens and that their features can predict metastatic outcome using an artificial neural network.

\section{Related Work}

The concept of EMT in cancer metastasis has been extensively studied in experimental models. \citet{mani2008mesenchyme} demonstrated that induction of EMT in mammary epithelial cells generates cells with stem cell properties. \citet{nieto2016emt} reviewed the evidence for EMT as a spectrum of hybrid states rather than a binary switch, with cells in intermediate states potentially having the highest metastatic capacity. \citet{pastushenko2018identification} identified distinct EMT states in skin tumours with differential metastatic potential, supporting the concept of a partial or hybrid EMT state as the most aggressive.

The role of CSC markers in oral cancer has been explored by several groups. CD24 and EpCAM have been used as markers to identify CSC populations in various carcinomas \citep{batlle2017cancer}. Vimentin expression in epithelial tumour cells has been associated with EMT and poor prognosis \citep{satelli2011vimentin}. However, the combination of these three markers to identify a disseminating EMT CSC population at the invasive front has not been previously attempted.

Machine learning approaches for predicting cancer outcomes from histopathological features have gained traction \citep{echle2021deep}, but applications specifically using quantified disseminating cell features are novel.

\section{Methods}

\subsection{Study cohort}
We analyzed 84 formalin-fixed, paraffin-embedded (FFPE) specimens from patients with primary oral squamous cell carcinoma. Cases were classified as metastatic (with regional lymph node metastasis) or non-metastatic based on pathological staging and clinical follow-up.

\subsection{Multi-marker immunofluorescence}
Triple immunofluorescence staining was performed using antibodies against EpCAM (epithelial marker), CD24 (stem cell marker), and Vimentin (mesenchymal marker). Secondary antibodies were conjugated to spectrally distinct fluorophores. DAPI was used for nuclear counterstaining. The staining protocol was optimized for FFPE oral cancer sections.

\subsection{Quantitative imaging analysis}
Over 12,000 imaging fields were captured across the 84 specimens, focusing on the invasive front. Automated image analysis pipelines identified single cells beyond the tumour body (disseminating cells) and quantified marker expression. Cells were classified by marker co-expression pattern: EpCAM$^+$ only, Vimentin$^+$ only, EpCAM$^+$/Vimentin$^+$ (double-positive), and EpCAM$^+$/CD24$^+$/Vimentin$^+$ (triple-positive).

\subsection{Artificial neural network}
A feedforward neural network was trained on features derived from disseminating cell counts and marker co-expression profiles to predict metastatic outcome. Performance was assessed using cross-validation on the full cohort of 84 specimens.

\section{Results}

\subsection{Identification of disseminating EMT CSCs at the invasive front}
Multi-marker immunofluorescence revealed distinct populations of cells at the invasive front of oral cancers. Triple-positive EpCAM$^+$/CD24$^+$/Vimentin$^+$ cells, representing a hybrid epithelial--mesenchymal phenotype with stem cell features, were identified as single cells disseminating beyond the tumour body into the surrounding stroma. These cells retained epithelial identity (EpCAM$^+$) while co-expressing the mesenchymal marker Vimentin, consistent with a partial EMT state.

\subsection{Marker co-expression analysis}
We systematically analyzed the association between marker co-expression patterns and dissemination/metastasis (Table~\ref{tab:markers}). EpCAM$^+$-only cells were predominantly non-disseminating and weakly associated with metastasis. Vimentin$^+$-only cells corresponded to mesenchymal stromal cells and were not predictive of metastasis. Double-positive EpCAM$^+$/Vimentin$^+$ cells showed moderate association with both dissemination and metastasis. The triple-positive EpCAM$^+$/CD24$^+$/Vimentin$^+$ population was the strongest predictor, being the phenotype most associated with active dissemination and metastatic outcome.

\begin{table}[h]
\centering
\caption{Marker co-expression patterns and their association with dissemination and metastasis.}
\label{tab:markers}
\begin{tabular}{lcc}
\toprule
Marker combination & Dissemination & Metastasis association \\
\midrule
EpCAM$^+$ only & Non-disseminating & Weak \\
Vimentin$^+$ only & Stromal & Not predictive \\
EpCAM$^+$/Vimentin$^+$ & Partial EMT & Moderate \\
EpCAM$^+$/CD24$^+$/Vimentin$^+$ & EMT CSC & Strong predictor \\
\bottomrule
\end{tabular}
\end{table}

\subsection{Enrichment in metastatic specimens}
Triple-positive disseminating cells were significantly enriched at the invasive front of metastatic compared to non-metastatic specimens ($p < 0.05$). The total count of single disseminating cells beyond the tumour body was also significantly higher in metastatic cases. This enrichment was specific to the triple-positive population; single-marker-positive cells did not show the same degree of enrichment.

\subsection{Neural network prediction of metastasis}
An artificial neural network trained on disseminating cell counts and marker co-expression features achieved a cross-validated accuracy of 87--89\% for predicting metastatic outcome across the 84 specimens (Table~\ref{tab:ann}). This performance demonstrates that features derived from disseminating EMT CSCs carry sufficient information to predict clinical metastasis with high accuracy.

\begin{table}[h]
\centering
\caption{Artificial neural network performance for metastasis prediction.}
\label{tab:ann}
\begin{tabular}{lc}
\toprule
Metric & Value \\
\midrule
Cross-validated accuracy & 87--89\% \\
Input features & Disseminating cell counts and marker co-expression \\
Model type & Artificial neural network \\
Validation & Cross-validation on 84 specimens \\
\bottomrule
\end{tabular}
\end{table}

\section{Discussion}

Our study provides the first direct observation of EMT cancer stem cells actively disseminating from human tumour tissue. By combining epithelial (EpCAM), stem cell (CD24), and mesenchymal (Vimentin) markers, we identified a triple-positive population at the invasive front of oral cancers that represents a hybrid EMT CSC state. The significant enrichment of these cells in metastatic specimens and the high accuracy of the neural network classifier trained on their features demonstrate that disseminating EMT CSCs have clinical predictive value.

The hybrid EMT state of these cells---co-expressing both epithelial and mesenchymal markers---is consistent with the emerging view that partial EMT, rather than complete mesenchymal conversion, underlies the most aggressive metastatic behavior \citep{pastushenko2018identification, nieto2016emt}. The retention of epithelial identity (EpCAM positivity) alongside mesenchymal features (Vimentin) suggests that these cells maintain the plasticity required for both dissemination and colonization at distant sites.

The addition of CD24 as a stem cell marker was critical for predictive accuracy. While double-positive EpCAM$^+$/Vimentin$^+$ cells showed moderate association with metastasis, the triple-positive population was a substantially stronger predictor. This supports the intertwined nature of EMT and stemness in driving metastatic capacity \citep{mani2008mesenchyme}.

The 87--89\% cross-validated accuracy of our neural network classifier compares favorably with existing prognostic markers for oral cancer metastasis. If validated in larger independent cohorts, this approach could be developed into a clinical tool for identifying patients at high risk of lymph node metastasis who might benefit from elective neck dissection or adjuvant therapy.

Limitations include the relatively small cohort size (84 specimens), the lack of an independent validation cohort, and the focus on a single cancer type. Additionally, the cross-sectional nature of histological analysis captures a snapshot of dissemination rather than the dynamic process itself.

\section{Conclusion}

We identify EpCAM$^+$/CD24$^+$/Vimentin$^+$ cells disseminating at the invasive front of human oral cancers as EMT cancer stem cells. These cells are significantly enriched in metastatic specimens, and their features predict metastasis with 87--89\% accuracy. This work bridges the gap between experimental models of EMT-driven metastasis and clinical observation, providing direct evidence that EMT CSCs disseminate from human tumours and that their presence is a clinically meaningful predictor of metastatic outcome.

\bibliographystyle{plainnat}
\bibliography{references}

\end{document}
